% ECONOMICS_PROBLEM: {"id": "C23-425", "title": "Continuous-treatment DiD", "jel_code": "C23", "source_id": "OP-0423"}
% !TeX program = xelatex
\documentclass[11pt,a4paper]{article}
\usepackage[margin=23mm]{geometry}
\usepackage{fontspec}
\setmainfont{Latin Modern Roman}
\usepackage{amsmath,amssymb,mathtools}
\usepackage{xcolor,enumitem,booktabs,longtable,array,xurl,hyperref}
\definecolor{muted}{HTML}{53636C}
\hypersetup{colorlinks=true,urlcolor=blue,linkcolor=blue}
\setlength{\parindent}{0pt}
\setlength{\parskip}{5pt}
\newcommand{\R}{\mathbb R}
\newcommand{\E}{\mathbb E}
\newcommand{\N}{\mathbb N}
\newcommand{\ind}{\mathbf 1}
\newcommand{\Var}{\operatorname{Var}}
\newcommand{\Cov}{\operatorname{Cov}}
\newcommand{\supp}{\operatorname{supp}}
\DeclareMathOperator*{\argmin}{arg\,min}
\DeclareMathOperator*{\argmax}{arg\,max}
\newcommand{\entrymeta}[3]{{\sffamily\small\color{muted}\textbf{\texttt{#1}}\quad|\quad #2\par #3\par}}
\newcommand{\Problemref}[1]{\texttt{#1}}
\newcommand{\JELref}[2]{\texttt{#2} (JEL \texttt{#1})}
\begin{document}
\section*{Continuous-treatment DiD}
\textbf{Problem ID:} \texttt{C23-425}\quad \textbf{JEL:} \texttt{C23}\par
% BEGIN ECONOMICS STATEMENT
\label{problem:OP-0423}
{\sffamily\small\color{muted}Original subject group: Econometrics, causal inference and measurement\par}
\label{definitions:problem:30007705}
\textbf{Audit classification:} Background; proposed extension. The paper is now AER forthcoming; both that publisher record and the NBER citation are verified. Continuous-treatment DiD is not generally unsolved, and the supplied selection restrictions define an editorial bounds exercise.
\subsection*{Definitions and precise research target}
Let units be observed before and after a policy, with baseline outcome \(Y_0\), nonnegative post-policy dose \(D\), and potential post-policy outcomes \(Y_1(d)\in[0,1]\) for feasible doses \(d\in[0,\bar d]\). Consistency gives observed \(Y_1=Y_1(D)\), and no anticipation makes \(Y_0\) unaffected by future dose. A conventional parallel-trends restriction compares \(E[Y_1(0)-Y_0\mid D=d]\) across dose groups; by itself it does not identify all comparisons between their treated potential outcomes. Fix a target dose interval \(B\) with positive probability and two doses \(d,d'\). Define
\[
\tau_B(d,d')=E[Y_1(d)-Y_1(d')\mid D\in B].
\]
Determine which additional selection restrictions, randomized dose information or bounded heterogeneity assumptions yield point identification or informative sharp bounds for this estimand from the observed panel law. For example, a declared Lipschitz constant \(L\) may restrict individual responses by \(|Y_1(d)-Y_1(d')|\leq L|d-d'|\), but its credibility requires external support. Distinguish dose-group average effects from causal derivatives, and state overlap and smoothness conditions whenever derivatives are targeted. The cited paper establishes several solutions under specified assumptions; this proposed restriction-and-bounds problem preserves residual selection concerns without describing continuous-treatment difference-in-differences as wholly unresolved.

\textbf{Definitions, assumptions and mathematical qualifications.} The observed law is that of \((Y_0,D,Y_1)\), with \(\Pr(D=0)>0\) when untreated units anchor trends. A conventional restriction is \(E[Y_1(0)-Y_0\mid D=d]=E[Y_1(0)-Y_0\mid D=0]\) for dose values almost everywhere on observed support. For continuous doses, conditional means at a particular point are defined only almost everywhere; pointwise and derivative targets additionally require continuity and overlap neighborhoods. The fixed subgroup \(B\) has \(\Pr(D\in B)>0\), and \(d,d'\) are interventions independent of the unit observed dose. Sharpness means that every value in the proposed set is generated by a joint potential-response law reproducing observations and all response restrictions.
\subsection*{Variants and assumption changes}
\begin{enumerate}
\item \textbf{Sourced solved benchmark.} Conventional parallel trends identifies dose-specific treatment-on-the-treated-type contrasts against zero, but comparisons across dose groups can mix response and selection; the paper analyzes stronger assumptions.
\item \textbf{Editorial response bounds.} Impose a validated Lipschitz bound, monotonic response and specified dose-selection restrictions; derive the identified set for \(\tau_B(d,d')\), distinguishing limits at exact doses from neighborhood-average effects.
\end{enumerate}
\subsection*{References and source audit}
\textbf{Checked source.} NBER Working Paper 32117 (2024), indexed abstract; AER forthcoming, DOI 10.1257/aer.20240137, publisher abstract. The paper supplies identification and estimation results and explains dose-selection limitations. Publisher primary abstract retrieved; NBER PDF returned HTTP 403. \url{https://www.nber.org/papers/w32117}.
\begin{enumerate}\raggedright
\item\hypertarget{definitions:source:30007705:1}{} Brantly Callaway; Andrew Goodman-Bacon; Pedro H. C. Sant'Anna. 2024. \emph{Difference-in-Differences with a Continuous Treatment}. NBER Working Paper 32117 (2024), indexed abstract; AER forthcoming, DOI 10.1257/aer.20240137, publisher abstract.
\url{https://www.nber.org/papers/w32117}.
DOI: \url{https://doi.org/10.3386/w32117}.
\end{enumerate}

\subsection*{Common conventions: Source mathematical conventions}

These conventions apply to each entry unless explicitly replaced.

\textbf{Models and identification.} A primitive model \(m\) specifies all probability laws, feasible choices, information, equilibrium and measurement rules used by its target. The admissible class \(\mathcal M\) is fixed independently of outcomes. If \(P_O(m)\) is its observable probability law and \(\tau(m)\) the target, its identified set at observable law \(P\) is
\[
 \mathcal I_\tau(P)=\{\tau(m):m\in\mathcal M,\ P_O(m)=P\}.
\]
Point identification means that this set is a singleton. An empty set signals incompatibility of data and assumptions. Coordinate bounds need not describe the joint set. Bounds are sharp only if every retained target is attainable or arbitrarily approachable by compatible models. Finite bounds require explicit restrictions on unobserved quantities; unrestricted confounding, hidden wealth, extrapolation or arbitrary equilibrium selection can yield uninformative sets.

\textbf{Causal interventions.} A potential outcome \(Y_i(p)\) is the outcome under a complete policy regime \(p\), including timing, financing, information and market spillovers. The target population and units are fixed before treatment. With interference, use the complete assignment vector or a justified exposure mapping; individual no-interference notation alone cannot identify an economy-wide intervention. A difference of mean potential outcomes does not require the cross-world joint distribution. Conditioning on post-treatment migration, survival, enrollment or employment changes the estimand and needs separate justification.

\textbf{Statistical statements.} An estimator, confidence set, forecast or test specifies a sampling law, model class and loss/coverage criterion. Uniform \((1-\alpha)\)-coverage means \(\inf_{m\in\mathcal M}\Pr_m\{\tau(m)\in C_n\}\geq1-\alpha\), or an explicitly asymptotic version. The whole parameter space trivially has coverage, so informativeness requires a stated width, risk or power objective. A forecasting improvement is relative to a named benchmark, information set and evaluation population.

\textbf{Equilibrium and optimization.} An equilibrium comprises feasible plans, optimality, beliefs where needed, budget constraints and all market/resource clearing. The primitives or a stated benchmark define each correspondence \(\mathcal E(p;m)\). Nonemptiness is assumed or proved, never inferred just from compact policy sets. A selected-equilibrium target uses a declared selection rule; a robust target quantifies over all permitted equilibria. Use suprema or infima unless attainment is established. An \(\varepsilon\)-optimal policy is feasible and within \(\varepsilon\) of the finite optimum value. Report an empty feasible set separately.

\textbf{Welfare and accounting.} Normative utility, interpersonal normalization, social weights, numeraire, discounting and the use of public receipts are primitives. Behavioral utility need not coincide with the welfare criterion. Household welfare includes ownership income and rents once; adding profits again to compensating variation generally double-counts them. A monetary equivalent is defined by an explicit transfer and price convention, with monotonicity and range conditions. Average effects, mechanism contributions, transfers and welfare losses are different targets. Infinite sums require convergence and dynamic feasible plans require terminal or no-Ponzi conditions.

\textbf{Source versions.} A working-paper date, revision date and journal publication year are distinguished. A DOI for a journal version is not automatically the DOI of an earlier manuscript. Locators refer to the stated version and printed pages where specified. A metadata-only check does not verify a claim about a full-text conclusion. Every entry's source subsection narrows the provenance claim accordingly.
% END ECONOMICS STATEMENT
\end{document}
